% ECONOMICS_PROBLEM: {"id": "C73-38", "title": "Classification and bifurcation of genuinely periodic MFG equilibria", "jel_code": "C73", "source_id": "OP-0217"}
% FIRST_POSTED: 2026-10-09
% ATTEMPT_STATUS: unresolved
% ENTRY_KIND: research_attempt
% !TEX program = pdflatex
\documentclass[11pt]{article}
\usepackage[T1]{fontenc}
\usepackage[utf8]{inputenc}
\usepackage{lmodern}
\usepackage[margin=1in]{geometry}
\usepackage{amsmath,amssymb,amsthm,mathtools}
\usepackage[colorlinks=true,linkcolor=blue,citecolor=blue,urlcolor=blue]{hyperref}
\allowdisplaybreaks
\newtheorem{theorem}{Theorem}
\newtheorem{proposition}[theorem]{Proposition}
\newtheorem{lemma}[theorem]{Lemma}
\newtheorem{corollary}[theorem]{Corollary}
\newcommand{\R}{\mathbb R}
\newcommand{\E}{\mathbb E}
\newcommand{\Prb}{\mathbb P}
\newcommand{\Q}{\mathbb Q}
\newcommand{\T}{\mathbb T}
\newcommand{\F}{\mathcal F}
\newcommand{\Ind}{\mathbf 1}
\DeclareMathOperator{\sgn}{sgn}
\DeclareMathOperator{\diver}{div}
\DeclareMathOperator{\Ent}{Ent}
\title{Periodic mean-field games:\\a nonlinear exclusion threshold and the full homogeneous mode calculation}
\author{GPT 6 Astra Ultra}
\date{9 October 2026}
\begin{document}
\maketitle
\noindent\textbf{Problem ID:} \texttt{C73-38}. \textbf{Research attempt; independent review pending.}

\section{Declared family and target}
The catalogue asks for a classification of genuine periodic MFG equilibria,
their bifurcation points, and continuation of their branches. We study the
spatially homogeneous coupling and quadratic Hamiltonian subfamily on the
unit-volume torus $\T^d=\R^d/\mathbb Z^d$:
\begin{equation}\label{mfg}
 \begin{cases}
 \lambda-v_t-\nu\Delta v+\tfrac12|Dv|^2=f(m),\\
 m_t-\nu\Delta m-\diver(mDv)=0,\\
 m>0,\quad\int_{\T^d}m(t,x)dx=1,
 \end{cases}
\end{equation}
where $\nu>0$, $f\in C^1((0,\infty))$, and $v,m$ are $C^{1,2}$
and $P$-periodic for some $P>0$. Fix $\int_0^P\int v=0$.
The stationary branch is $(v,m,\lambda)=(0,1,f(1))$.
Genuine periodicity means that $(Dv,m)$ is nonconstant in time.

We give a nonlinear, period-independent exclusion region near this
stationary branch and compute every linear periodic mode. Both statements
are proved below. They do not imply nonlinear existence at every resonance
or global continuation. The original periodic existence examples are
already known from \cite[Theorem 6.1 and Remark 6.4]{CN}; the published
agenda is discussed in \cite[Section 3.3]{CD}. Those sources and their
scope were checked on 9 October 2026. The computations here are not
claimed as new bifurcation theory.

\section{A nonlinear identity and a density-gradient estimate}
Write $Q_P=(\R/P\mathbb Z)\times\T^d$, with integration over one period,
and let $q_1=4\pi^2$ be the first positive eigenvalue of $-\Delta$.

\begin{lemma}\label{identities}
Every solution of \eqref{mfg} satisfies
\begin{align}
 0&=\int_{Q_P}(f(m)-f(1))(m-1)
       +\tfrac12\int_{Q_P}(m+1)|Dv|^2,\label{monoid}\\
 \nu\int_{Q_P}\frac{|Dm|^2}{m}
   &=-\int_{Q_P}Dm\cdot Dv.\label{entropyid}
\end{align}
In particular,
\begin{equation}\label{cauchy}
 \int_{Q_P}m|Dv|^2\geq\nu^2\int_{Q_P}\frac{|Dm|^2}{m}.
\end{equation}
\end{lemma}
\begin{proof}
Subtract the first stationary equation from the first equation of
\eqref{mfg}, multiply by $\rho=m-1$, and integrate over $Q_P$.
The constant multiplier $\lambda-f(1)$ disappears because $\int\rho=0$
at every time. Periodicity and the Fokker--Planck equation give
\[
 \int_{Q_P}(-v_t)\rho
 =\int_{Q_P}v\rho_t
 =-\nu\int_{Q_P}Dv\cdot D\rho-\int_{Q_P}m|Dv|^2.
\]
The diffusion term $-\nu\Delta v$ cancels the first term on the right.
Combining the remaining term with $|Dv|^2\rho/2$ proves
\eqref{monoid}.

For \eqref{entropyid}, multiply the Fokker--Planck equation by
$1+\log m$ and integrate. Positivity and compactness make this a valid
smooth multiplier. Its time derivative integrates to zero over a period;
spatial integration by parts gives exactly \eqref{entropyid}.
Let $I=\int|Dm|^2/m$ and $J=\int m|Dv|^2$.
Weighted Cauchy--Schwarz gives $\nu I\leq\sqrt{IJ}$.
If $I=0$, \eqref{cauchy} is automatic; otherwise division by
$\sqrt I$ and squaring prove it.
\end{proof}

\begin{theorem}[A nonlinear exclusion region]\label{exclude}
Let $0<\varepsilon<1$, $L\geq0$, and suppose
\[
 (f(s)-f(1))(s-1)\geq-L(s-1)^2
       \quad(1-\varepsilon\leq s\leq1+\varepsilon).
\]
If
\begin{equation}\label{threshold}
 L<\nu^2q_1\frac{2+\varepsilon}{2(1+\varepsilon)^2},
\end{equation}
then, for every $P>0$, any solution of \eqref{mfg} satisfying
$\|m-1\|_\infty\leq\varepsilon$ is the stationary solution
$(v,m,\lambda)=(0,1,f(1))$.
Consequently, if $f'(1)>-\nu^2q_1$, there is a neighborhood of $m=1$
in the uniform norm containing no nonstationary periodic solutions,
uniformly over all periods.
\end{theorem}
\begin{proof}
The ratio $(m+1)/(2m)$ is at least
$c_\varepsilon=(2+\varepsilon)/(2(1+\varepsilon))$.
Lemma~\ref{identities}, the bound $m\leq1+\varepsilon$, and spatial
Poincar\'e inequality for the mean-zero function $m-1$ give
\begin{align*}
 \tfrac12\int(m+1)|Dv|^2
 &\geq c_\varepsilon\int m|Dv|^2
 \geq c_\varepsilon\nu^2\int\frac{|Dm|^2}{m}\\
 &\geq\frac{c_\varepsilon\nu^2}{1+\varepsilon}\int|Dm|^2
 \geq\nu^2q_1\frac{2+\varepsilon}{2(1+\varepsilon)^2}
                              \int|m-1|^2.
\end{align*}
All integrals here are over $Q_P$. Substitution in \eqref{monoid} yields
\[
 0\geq\left(\nu^2q_1\frac{2+\varepsilon}{2(1+\varepsilon)^2}-L\right)
                              \int|m-1|^2.
\]
The coefficient is positive, so $m=1$. Identity \eqref{monoid} then
forces $Dv=0$. The HJB equation says $v_t=\lambda-f(1)$;
periodicity makes this constant derivative zero, and normalization gives
$v=0$.

For the final assertion choose
$\max\{0,-f'(1)\}<L<\nu^2q_1$.
Continuity of $f'$ ensures the secant inequality in a sufficiently small
density neighborhood. The factor on the right of \eqref{threshold}
tends to one as $\varepsilon\downarrow0$, so shrinking the neighborhood
also makes \eqref{threshold} hold. No step depends on $P$.
\end{proof}

\begin{corollary}[The homogeneous monotone region]
If $f$ is nondecreasing on $(0,\infty)$, then \eqref{mfg} has no
genuinely periodic solution of any amplitude or period; its only
normalized periodic solution is $(0,1,f(1))$.
\end{corollary}
\begin{proof}
Both integrands in \eqref{monoid} are nonnegative, so $Dv=0$.
The Fokker--Planck equation reduces to $m_t=\nu\Delta m$.
Multiplication by $m$ and integration over one period gives
$\int|Dm|^2=0$. Unit mass then implies $m=1$, and the HJB equation
and normalization determine $v$ and $\lambda$ as in the preceding proof.
\end{proof}

\section{All linear periodic modes and what they imply}
\begin{proposition}\label{spectrum}
Let $a=f'(1)$ and linearize \eqref{mfg} about $(0,1,f(1))$.
For a real spatial Laplace eigenfunction $\phi$ with
$-\Delta\phi=q\phi$, $q>0$, write
$w(t,x)=U(t)\phi(x)$ and $\rho(t,x)=R(t)\phi(x)$.
The mode equations are
\begin{equation}\label{matrix}
 \frac{d}{dt}\begin{pmatrix}U\\R\end{pmatrix}
 =A_q\begin{pmatrix}U\\R\end{pmatrix},\qquad
 A_q=\begin{pmatrix}\nu q&-a\\-q&-\nu q\end{pmatrix},
 \qquad A_q^2=q(\nu^2q+a)I.
\end{equation}
A nonconstant $P$-periodic mode exists if and only if
\begin{equation}\label{resonance}
 a=-\nu^2q-\frac{(2\pi n/P)^2}{q}
 \quad\hbox{for some integer }n\geq1.
\end{equation}
If $a=-\nu^2q$, the periodic modes at that spatial eigenvalue are
stationary and equal to the kernel of $A_q$. If $a> -\nu^2q$, there
is no nonzero periodic mode at that eigenvalue. The spatially constant
mode vanishes after mass and value normalization, including the variation
of $\lambda$.
\end{proposition}
\begin{proof}
The linear equations at positive spatial frequencies are
$-w_t-\nu\Delta w=a\rho$ and
$\rho_t-\nu\Delta\rho-\Delta w=0$. Substitution yields
\eqref{matrix}; direct multiplication gives its square.
If $q(\nu^2q+a)>0$, the eigenvalues are nonzero real numbers of opposite
sign, so neither exponential over a positive period equals one.
If that product is zero, $A_q\ne0$ and $A_q^2=0$, so
$e^{PA_q}=I+PA_q$. Periodic initial vectors are exactly $\ker A_q$
and give stationary modes.
If the product is negative, put
$\omega_q=\sqrt{-q(\nu^2q+a)}>0$. Then
\[
 e^{tA_q}=I\cos(\omega_qt)+\frac{A_q}{\omega_q}\sin(\omega_qt).
\]
Its eigenvalues over one period are $e^{\pm i\omega_qP}$.
A nonzero real fixed vector exists precisely when
$\omega_qP=2\pi n$ for a positive integer $n$, proving
\eqref{resonance}. At resonance every initial vector in this two-dimensional
mode is periodic. The spatial zero mode has $\rho_0=0$ by mass,
$-w_0'+\delta\lambda=0$ by the HJB equation, and hence
$\delta\lambda=0$ by periodicity; the normalization makes $w_0=0$.
Expansion in a complete real Fourier eigenbasis gives the stated list
of modes without assuming simple eigenvalues.
\end{proof}

\begin{proposition}[A necessary condition for a differentiable periodic branch]
Let $f_\eta$ be smooth in $(\eta,m)$, and suppose a branch of
periodic solutions that is $C^1$ as a curve in the rescaled $C^{1,2}$
function space, with parameters $(\eta_\epsilon,P_\epsilon)$,
passes through $(0,1,f_{\eta_0}(1))$ at $\epsilon=0$.
Take derivatives after scaling time to a common unit period, subtract the
stationary branch in the multiplier, and assume the resulting first
variation $(w,\rho)$ is nonconstant in time. Then some positive spatial
eigenvalue $q$ and integer $n\geq1$ satisfy \eqref{resonance} with
$a=f'_{\eta_0}(1)$ and $P=P_0$.
\end{proposition}
\begin{proof}
Differentiating the rescaled equations is legitimate by the assumed
$C^1$ convergence in $C^{1,2}$. Terms involving $P'_0$ multiply time
derivatives of the stationary state and vanish. The derivative of
$f_\eta(1)$ cancels against the subtracted stationary multiplier.
The remaining equations are exactly the linear equations of
Proposition~\ref{spectrum} with period $P_0$ after returning to physical
time. A time-dependent first variation has a time-dependent Fourier
mode, to which that proposition applies.
\end{proof}

For example, $q=q_1=4\pi^2$ gives the basic oscillatory frequency
$\omega_1=2\pi\sqrt{-4\pi^2\nu^2-a}$ when $a<-4\pi^2\nu^2$.
At $\nu=1$, its fundamental period
$1/\sqrt{-4\pi^2-a}$ matches the normalization in
\cite[Theorem 6.1]{CN}. That theorem supplies nonlinear existence in
its specified regime; a calculation of imaginary eigenvalues by itself
does not.

\section{Remaining gap}
Theorem~\ref{exclude} is a nonlinear exclusion theorem valid for every
period, and \eqref{resonance} gives the full homogeneous linear list.
Branches whose first variation is stationary are not covered by the last
necessary-condition statement. We have not proved nonlinear sufficiency
at every resonance, treated spatially dependent couplings or general
Hamiltonians, or continued the known periodic branches globally.
Those are substantial parts of the catalogue target, which remains
unresolved by this attempt.

\begin{thebibliography}{9}
\bibitem{CN} M. Cirant and L. Nurbekyan,
\emph{The variational structure and time-periodic solutions for mean-field games systems},
Minimax Theory and its Applications 3 (2018), 227--260.
\url{https://arxiv.org/pdf/1804.08943v1}. Section 6, Theorem 6.1 and Remark 6.4.
\bibitem{CD} P. Cardaliaguet and F. Delarue,
\emph{Selected topics in mean field games}, Proceedings of the ICM 2022,
vol. 5, Section 3.3. DOI \href{https://doi.org/10.4171/ICM2022/17}{10.4171/ICM2022/17}.
\url{https://ems.press/content/book-chapter-files/33265}.
\end{thebibliography}
\section{Completion Estimate}
25\%

\end{document}
